% citum-example.tex — Full demonstration of the Citum LuaLaTeX package
%
% Usage:
%   lualatex citum-example.tex
%
% Requirements:
%   1. The 'citum' engine must be compiled with the 'ffi' feature.
%   2. Either set CITUM_LIB_PATH to the full path of libcitum_processor.dylib/.so,
%      or place the library in the same directory as this file.
%   3. Either set CITUM_LUA_PATH to the full path of citum.lua,
%      or place citum.lua and citum.sty in the same directory as this file.

\documentclass{article}
\usepackage[margin=1.2in]{geometry}
\usepackage{fontspec}
\usepackage{hyperref}

% Citum citation package
%   style   = path to a Citum YAML style file (extension optional)
%   bibfile = path to a Citum YAML bibliography (extension optional)
%   locale  = optional locale name (e.g. fr-FR)
\usepackage[
    style   = apa-7th,
    bibfile = example-refs,
    % locale  = en-US
]{citum}

\title{Modern Citations in Lua\LaTeX\\
       A Demonstration of Citum in Lua\LaTeX}
\author{Citum Labs}
\date{March 2026}

\begin{document}

\maketitle

\section{Introduction}

This document demonstrates the Citum citation engine 
integrated directly into Lua\LaTeX. 

Citations are rendered live by the Citum Rust processor via a LuaJIT FFI 
binding. There is no need for external tools like Biber or BibTeX; the 
rendering happens during the Lua\LaTeX{} pass itself.

\section{Citation Examples}

\subsection{Standard Citations}

You can use the standard \verb|\cite| command for non-integral citations. 
For example, we might discuss the structure of scientific revolutions 
\cite{kuhn1962}.

Optional locators are supported: \cite[p. 52]{kuhn1962}.

\subsection{Integral Citations}

Integral citations (where the author is part of the sentence) use 
\verb|\textcite|. As \textcite{lecun2015} argues, deep learning is a 
powerful tool. 

Locators work here too: \textcite[ch. 3]{lecun2015} provides more detail.

\subsection{Multiple Citations}

Multiple keys can be cited at once using \verb|\cites{key1, key2}|: 
\cites{kuhn1962, lecun2015}.

\subsection{Complex Multi-item Citations}

For per-item locators or prefixes in a group, use the \verb|\citestart|, 
\verb|\citeitem|, and \verb|\citeend| block:
\citestart
  \citeitem[see]{kuhn1962}
  \citeitem[also]{lecun2015}
\citeend

\section{Technical Advantages}

The Citum processor renders these citations --- non-integral parentheticals, 
integral name-date combinations, and the final bibliography --- using a 
unified rendering engine. This ensures that disambiguation, sorting, and 
formatting are consistent across all output formats.

\section{Bibliography}

The bibliography is printed with \verb|\printbibliography| (or 
\verb|\printcitumbibliography|).

\printbibliography

\end{document}
